% Generated draft: requires mathematical and source review.
\documentclass[11pt]{article}
\usepackage[margin=25mm]{geometry}
\usepackage{fontspec}
\setmainfont{Latin Modern Roman}
\usepackage{amsmath,amssymb,mathtools,unicode-math}
\setmathfont{Latin Modern Math}
\usepackage{enumitem,needspace,xurl,hyperref}
\setlength{\parindent}{0pt}
\setlength{\parskip}{5pt}
\setlength{\emergencystretch}{3em}
\newcommand{\N}{\mathbb{N}}
\newcommand{\Z}{\mathbb{Z}}
\newcommand{\Q}{\mathbb{Q}}
\newcommand{\R}{\mathbb{R}}
\newcommand{\C}{\mathbb{C}}
\newcommand{\F}{\mathbb{F}}
\newcommand{\A}{\mathbb{A}}
\newcommand{\PP}{\mathbb{P}}
\newcommand{\E}{\mathbb{E}}
\newcommand{\eps}{\varepsilon}
\newcommand{\dd}{\,\mathrm{d}}
\begin{document}
\section*{UnsolvedMath ID: 800774}
\textbf{Source problem code:} catalogue entry
\subsection{Definitions and mathematical statement}
Let \(X\) be a compact K\"ahler manifold whose holomorphic tangent bundle \(T_X\) is pseudo-effective, in the sense intended by the source. The supplied record does not specify the definition of pseudo-effectivity for a vector bundle; this hypothesis requires clarification before the statement can be considered fully precise.

A finite \emph{\'etale covering} \(p:Y\to X\) means a finite surjective unramified holomorphic map. Let
\[
a_Y:Y\longrightarrow A_Y:=\operatorname{Alb}(Y)
\]
be the Albanese map, namely the universal holomorphic map from \(Y\) to a complex torus after choosing a base point.

The source uses \(F\) without defining it. The following formulation provisionally interprets \(F\) as a smooth general fiber of \(a_Y\) over its image. The record does not supply surjectivity, connectedness of fibers, or smoothness of the Albanese map as hypotheses or established conclusions. Consequently, this interpretation of \(F\), including whether a whole fiber or a connected component is intended, requires review.

For a smooth compact complex manifold \(F\) of dimension \(d\), write
\[
K_F=\bigwedge^d\Omega_F^1,
\qquad -K_F:=K_F^{-1}.
\]
Here bigness of \(-K_F\) means
\[
\limsup_{m\to\infty}
\frac{h^0(F,K_F^{-m})}{m^d}>0,
\]
and projectivity means that \(F\) admits a holomorphic embedding into some complex projective space.

A maximal rationally connected (MRC) fibration is understood in its usual birational sense: its general fibers are rationally connected and, through a very general point of the total space, the fiber represents the equivalence class generated by chains of rational curves. Its base is determined up to birational equivalence. Rational connectedness means that two general points can be joined by a rational curve.

\textbf{Questions attributed to Matsumura (2022), Problem 9.} For every \(X\) satisfying the stated hypothesis, after replacing \(X\) by an appropriate finite \'etale covering \(Y\), ask:
\begin{enumerate}
\item Can one prove that \(-K_F\) is big, or that \(F\) is projective?
\item Is \(a_Y:Y\to A_Y\) an MRC fibration of \(Y\)?
\end{enumerate}
The phrase ``appropriate finite \'etale covering'' is read existentially. The source does not specify whether the two questions must hold for the same covering. In the first question, it also does not specify whether bigness is the primary target with projectivity as a weaker alternative, or whether only the displayed disjunction is requested. In the second question, taking the entire Albanese torus as the MRC base would require the Albanese map to be surjective; replacing that target by its image is a distinct formulation that is not silently assumed here.
\subsection{Short English statement}
Matsumura asks whether a compact K\"ahler manifold with pseudo-effective tangent bundle has, after a suitable finite \'etale cover, Albanese fibers with big anti-canonical bundle or at least projective fibers, and whether its Albanese map gives the maximal rationally connected fibration. The meanings of pseudo-effectivity and of the undefined fiber \(F\), together with the precise covering and target conventions, require clarification.
\subsection{Sources}
\begin{itemize}
\item Local UnsolvedMath registry: permanent ID \(800774\), exact source problem code \texttt{catalogue entry}, title ``Positively Curved Projective Varieties --- Problem 9.'' The supplied record attributes the problem to Matsumura, with proposed year 2022, and labels it \emph{open}; this is a supplied status label, not an independently verified mathematical assessment.
\item Supplied source attribution: Matsumura, \emph{Open Problems on Structure of Positively Curved Projective Varieties} (2022), Problem 9. Supplied URL: \url{https://arxiv.org/abs/2211.09153}. The background records extraction by \texttt{arxiv\_latex} and access on \texttt{2026-10-16}; these are supplied provenance statements, not access or verification performed for this draft.
\item Registry provenance: upstream ID \(800774\), revision \texttt{d371ee917ac9647fbb06cd6fc529deec6f3dda82}; supplied creation and update timestamps both \texttt{2026-10-16T14:00:00Z}. No further literature references or independent source verification are supplied.
\end{itemize}
\end{document}
